%@ pages 54-60
%@ status translated
% Original section 5, "La structure de type Bousfield-Kan d'après Hirschhorn", pp. 54--60.
%
% Translation notes (the mathematics is left exactly as printed):
%  * The authors write "n-précat de Segal" (short for "n-précatégorie"),
%    the objects of the category nSePC (defined in \S2).  It is rendered
%    "Segal n-precategory", to be kept distinct from the "n-préchamp de
%    Segal" = "Segal n-prestack" (a presheaf with values in nSePC), which
%    is the other notion used in this section.
%  * nSePC and nSePCh are both printed in the source (nSePCh in the
%    homotopy category Ho(nSePCh(X^gro)) and in the Gabriel-Zisman
%    sentence); both spellings are kept.
%  * "libre addition de cellules" is rendered "free addition of cells",
%    "librement engendrée" as "freely generated".

\clearpage
\section{The Bousfield--Kan-type structure after Hirschhorn}
\label{bousfieldkanpage}

We wish to indicate here another way of obtaining a closed model category
for Segal $n$-stacks, based on Hirschhorn's book \cite{Hirschhorn} and
originating in a construction of Bousfield--Kan \cite{BousfieldKan} (which in
turn goes back to Quillen \cite{Quillen}) for the case of simplicial
presheaves.

In \cite{Hirschhorn}, Hirschhorn first defines the notion of a \emph{closed
model category generated by cofibrations}. We have taken this notion up in
statement \ref{dhklemme} above, which can therefore serve as a definition.
Next, Hirschhorn shows that, for every small category $\Xx$ and every cmc $M$
generated by cofibrations, the category of presheaves on $\Xx$ with values in
$M$ admits a structure of closed model category denoted
$$
M^{\Xx ^o}.
$$
For $M=EnsSpl$ the construction is due to Bousfield--Kan (\cite{BousfieldKan},
p.~314). Their construction was in turn a direct application of Quillen's
construction \cite{Quillen} of a cmc for simplicial objects in a suitable
category, in this case that of presheaves of sets. For this reason we shall
call \emph{HBKQ structure (i.e.\ Hirschhorn--Bousfield--Kan--Quillen)} this
cmc structure defined by Hirschhorn on $M^{\Xx ^o}$.

For $M=nSePC$, the category of Segal $n$-precategories, the HBKQ structure is
different from that of Theorem \ref{cmf}, our model category of Segal
$n$-prestacks for the coarse topology. The underlying categories are the
same, as are the weak equivalences (which are the object-by-object weak
equivalences). On the other hand, the HBKQ fibrations are the morphisms
$f:A\rightarrow B$ inducing for every $X$ a fibration
$f(X): A(X)\rightarrow B(X)$ ---this is therefore a larger class than that of
the fibrations of Theorem \ref{cmf}. And whereas we take all injections as
cofibrations, the class of HBKQ cofibrations is smaller: they are the
retractions of morphisms which are obtained by a (transfinite) sequence of
free additions of cells, where a \emph{free addition of cells} is a
cofibration $A\rightarrow A'$ freely generated by a cofibration
$A(X)\rightarrow C'$ in $M$ (for a single object $X$). More precisely, if one
denotes by $p_X: \Xx /X\rightarrow \Xx$ the natural functor, and by
$\underline{C}$ the constant diagram on $\Xx /X$ with values in an object $C$
of $M$, then a free addition of cells is a coproduct of the form
$$
A' = A \cup ^{p_{X,!}(C)}p_{X,!}(C')
$$
for a cofibration $C\rightarrow C'$ in $M$ (in the preceding expression we
had $C=A(X)$).

The interest of the morphisms obtained by free addition of cells lies in the
following fact: if $A\rightarrow A'$ is obtained by adding a cell
$A(X)\rightarrow C'$ over $X$, and if $B$ is another diagram, then the
extension of a morphism $f:A\rightarrow B$ to $f':A'\rightarrow B$ is nothing
other than the extension of the morphism $f(X): A(X)\rightarrow B(X)$ to a
morphism $C'\rightarrow B(X)$.

Since the two categories and their subcategories of weak equivalences are the
same, the two homotopy theories are the same (by the Dwyer--Kan localisation
theory \cite{DwyerKan1} \cite{DwyerKan2} \cite{DwyerKan3}).

{\em Conclusion:} One could have taken, as a closed model category for the
Segal $n$-prestacks on $\Xx$ with the coarse topology, the HBKQ structure on
$M^{\Xx ^o}$ with $M$ the closed model category $nSePC$ of Segal
$n$-precategories (Theorem \ref{SeCmf}). In the HBKQ structure the fibrations
are simply the morphisms which are object-by-object fibrations of Segal
$n$-precategories.

{\em Remark:} If $\Xx$ admits fibre products then the HBKQ cofibrations are
stable under restriction to the sites $\Xx /X$ and also stable under direct
product. This allows one to apply the theory of internal $Hom$'s which will
be explained in \S11 below.

To legitimise this other possible point of view, it remains to see how to
integrate a Grothendieck topology into this theory. Let again $M$ be the
category of Segal $n$-precategories of Theorem \ref{SeCmf}. On $M^{\Xx ^o}$,
which is the category of Segal $n$-prestacks, one keeps the notion of HBKQ
cofibration; one uses the notion of $\Gg$-weak equivalence defined above; and
one defines the \emph{$\Gg$-fibrations of HBKQ} as the morphisms which
possess the lifting property with respect to the HBKQ cofibrations which are
$\Gg$-weak equivalences.

\begin{theorem}
\label{cmfHirschho}
The category of Segal $n$-prestacks equipped with the HBKQ cofibrations, the
$\Gg$-weak equivalences, and the $\Gg$-fibrations of HBKQ, is a closed model
category.
\end{theorem}
{\em Proof:}
A morphism is a $\Gg$-trivial cofibration of HBKQ if and only if it is a
cofibration of HBKQ and also a $\Gg$-trivial cofibration for the structure
\ref{cmf}. Hirschhorn proves that his cofibrations (for the coarse topology)
are stable under coproduct \cite{Hirschhorn}. Hence the preservation of
$\Gg$-trivial HBKQ cofibrations under coproduct is an immediate consequence
of the same result for the structure of Theorem \ref{cmf}.
The other properties are also immediate.
\eop


The reader may find this closed model category structure more congenial than
that of \ref{cmf}. We have kept the structure of \ref{cmf} as the basic
notion because it is compatible with Jardine's structure (which has become
standard, notably in $K$-theory) for simplicial presheaves, i.e. for the case
of Segal $0$-prestacks.

\subnumero{Computing homotopy classes of maps}

The HBKQ structure makes it possible to compute homotopy classes of maps more
easily. Let us fix a category $\Xx$ equipped with the coarse topology. The
homotopy category $Ho(nSePCh(\Xx ^{\rm gro}))$ is defined as the
Gabriel--Zisman localisation of $nSePCh(\Xx )$ by the subcategory of
object-by-object weak equivalences. In particular
$Ho(nSePCh(\Xx ^{\rm gro}))$ does not depend on the choice of cmc structure
compatible with these weak equivalences; in particular one may use the HBKQ
structure to compute the sets $Hom$ in the homotopy category, i.e. the sets
of homotopy classes of maps.

Let us begin with a remark of a general nature: if $E$ is a Segal
$n$-precategory and if $\underline{E}$ is the constant Segal $n$-prestack on
$\Xx$ with values $E$, then the direct product with $\underline{E}$ respects
the HBKQ cofibrations (with no hypothesis on the existence of fibre products
in $\Xx$, and without assuming that $\underline{E}$ is itself cofibrant).
Indeed, if $A\rightarrow A'$ is obtained by adding a cell $A(X)\rightarrow C$
over the object $X$, then $A\times \underline{E}\rightarrow A'\times
\underline{E}$ is obtained by adding the cell $A(X)\times E \rightarrow C\times
E$. More generally, if $E\rightarrow E'$ is a cofibration of Segal
$n$-precategories and $A\rightarrow A'$ is an HBKQ cofibration, then
$$
A\times \underline{E}' \cup ^{A\times \underline{E}}A'\times \underline{E}
\rightarrow A'\times \underline{E}'
$$
is again an HBKQ cofibration.

Let $A$ and $B$ be two Segal $n$-prestacks on $\Xx$ and suppose that $B(X)$ is
fibrant for every $X\in \Xx$. Then $B$ is a fibrant object for the HBKQ
structure (for the coarse topology). If one chooses an HBKQ-cofibrant
replacement $C\rightarrow A$ (object-by-object weak equivalence with $C$
HBKQ-cofibrant) then every element of $[A,B]_{\Xx ^{\rm gro}}$ comes from a
morphism $C\rightarrow B$. Let us denote by $\overline{I}$ the category with
two objects $0,1$ and an isomorphism between them (regarded as a Segal
$n$-category, see ``Induction'' in \S2). And let us also denote by
$\overline{I}$ the constant Segal $n$-prestack with values $\overline{I}$. The
diagram
$$
C\times \{ 0,1\} \rightarrow C\times \overline{I}\rightarrow C
$$
constitutes a ``$C\times I$ object'' in Quillen's terminology \cite{Quillen}
for the HBKQ structure, and in particular the first map is a cofibration.
Hence two maps $C\rightarrow B$ are homotopic, i.e. induce the same element of
$[A,B]_{\Xx ^{\rm gro}}$, if and only if they are connected by a homotopy of
the form $C\times \overline{I}\rightarrow B$. We have thus obtained the
following statement.

\begin{lemma}
\label{howtocalc}
Let $A$ be a Segal $n$-prestack on $\Xx$ and fix a replacement $C\rightarrow
A$ with $C$ HBKQ-cofibrant. Then for every Segal $n$-prestack $B$ such that
$B(X)$ is a fibrant Segal $n$-category for every $X\in \Xx$, the set of
homotopy classes of maps
$$
[A,B]_{\Xx ^{\rm gro}}
$$
is computed as the set of morphisms of Segal $n$-prestacks $C\rightarrow B$,
modulo the equivalence relation which identifies two morphisms if they are
connected by a homotopy of the form $C\times \overline{I}\rightarrow B$.
\end{lemma}
\eop


\subnumero{A uniform expression for descent data}

One can use this lemma \ref{howtocalc} to give a uniform expression for
descent data with values in a Segal $n$-prestack $A$ such that all the $A(X)$
are fibrant Segal $n$-categories. For this discussion one needs $n\geq 1$ (in
the case $n=0$, one must replace $A$ by $\Pi _{1,Se}\circ A$, see \S10 for
example).

Let $X\in \Xx$ and let $\Bb \subset \Xx /X$ be a sieve covering $X$. Let us
choose an HBKQ-cofibrant replacement
$$
D\rightarrow \ast _{\Bb}
$$
(let us say, to fix ideas, that we are speaking here of Segal $n$-prestacks
on $\Bb$---but one may also take the $p_!$ to find oneself on $\Xx /X$ or on
$\Xx$). Now, by lemma \ref{howtocalc}, we know that for every Segal
$n$-prestack $A$ such that all the $A(X)$ are fibrant, the set of descent data
$$
[\ast _{\Bb}, A|_{\Bb}]_{\Bb ^{\rm gro}}
$$
is expressed as the set of morphisms $D\rightarrow A|_{\Bb}$ modulo the
relation which identifies two morphisms if they are connected by a homotopy
$D\times \overline{I}\rightarrow A|_{\Bb}$.

Suppose now that the site $\Xx$ admits fibre products and that the sieve
$\Bb$ comes from a covering family $\Uu = \{ U_{\alpha }\rightarrow X\}$. In
this case one can give an explicit formula for an HBKQ-cofibrant replacement
$D\rightarrow \ast _{\Bb}$.

For this, let $\overline{I}^{(p)}$ be the category with $p+1$ objects denoted
$0',\ldots , p'$ and an isomorphism between every pair of objects (it is in
some sense the ``symmetric product'' of $p$ copies of $\overline{I}$). Let us
regard it as a Segal $n$-precategory (by ``induction'', see \S2). One obtains
a functor denoted
$$
R: \Delta \rightarrow nSePC,\;\;\;\;
R(p):= \overline{I}^{(p)}.
$$
This functor equips $nSePC$ with a simplicial cmc structure \cite{Quillen}. In
particular, if $K$ is a simplicial set (in this case one does not need $K$ to
be finite) and $E$ a Segal $n$-precategory, one obtains the Segal
$n$-precategory $K\otimes ^RE$; in our case one can write
$$
K\otimes ^R E = (K\otimes ^R\ast )\times E.
$$
Here $K\otimes ^R\ast$ is just the coproduct of the $R(p)$ over the
$p$-simplices of $K$ (taken according to the usual procedure for realising a
simplicial set). We shall say that $K\otimes ^R\ast$ is the \emph{realisation
of $K$ with respect to $R$}.

This construction being functorial, it extends to presheaves of simplicial
sets: if $K$ is a presheaf of simplicial sets on $\Xx$ one obtains the Segal
$n$-prestack $K\otimes ^R\ast$.

Now our covering family $\Uu$ represents a presheaf of sets: it is simply the
presheaf disjoint union of those represented by the $U_{\alpha}$. One obtains
the simplicial presheaf
$$
P(\Uu ):= \ldots \Uu \times _X \Uu \tworightarrows \Uu
$$
i.e. $P(\Uu )_k=\Uu \times _X \ldots \times _X\Uu $ ($k+1$ times). This is the
\emph{nerve of the covering $\Uu$}. Applying the construction ``realisation
with respect to $R$'' one obtains the Segal $n$-prestack
$$
D:= P(\Uu )\otimes ^R \ast .
$$
The morphism $D\rightarrow \ast _{\Bb}$ is an object-by-object equivalence.
Indeed $P(\Uu )(Y)$ is contractible if $Y\in \Bb$ and empty otherwise. Observe
that the object $D$ is defined as an $n$-prestack on $\Xx /X$, so that to
define the morphism $D\rightarrow \ast _{\Bb}$ one must interpret
$\ast _{\Bb}$ as an $n$-prestack on $\Xx /X$.

We claim that $D$ is HBKQ-cofibrant. In fact it is interesting to see
explicitly how to obtain $D$ from $\emptyset$ by a sequence of free additions
of cells. First one adds $R(0)=\ast$ over each element $U_{\alpha}$ of the
covering family. Let us call $D^0$ the result obtained. Next one adds
$R(1)=\overline{I}$ over each $U_{\alpha} \times _XU_{\beta}$, via the diagram
$$
D^0(U_{\alpha} \times _XU_{\beta}) \leftarrow p_1^{\ast}R(0)\sqcup
p_2^{\ast}R(0) = \{ 0,1\} \hookrightarrow \overline{I}.
$$
Let us denote by $D^1$ the result obtained. More generally, let us denote by
$\partial R(p)$ the ``boundary'' of $R(p)$, i.e. the union in $R(p)$ of the
$p+1$ faces of the form $R(p-1)$. One adds $R(2)$ over each
$U_{\alpha}\times _XU_{\beta} \times _XU_{\gamma}$ via the diagram
$$
D^1(U_{\alpha}\times _XU_{\beta}
\times _XU_{\gamma})\leftarrow \partial R(2) \hookrightarrow  R(2)
$$
to obtain $D^2$. At each step one adds to $D^{p-1}$ a copy of $R(p)$ over each
$$
U_{\alpha _0,\ldots , \alpha _p}:=U_{\alpha _0}
\times _X\ldots \times _X U_{\alpha_p},
$$
via the diagram
$$
D^{p-1}(U_{\alpha _0} \times _X\ldots \times _X
U_{\alpha_p})\leftarrow \partial R(p) \hookrightarrow  R(p).
$$
Finally $D$ is the colimit (increasing union) of the $D^p$.

One notes that $D$ can be regarded as a Segal $n$-prestack over $\Bb$ and the
description above shows that as such (as well as a Segal $n$-prestack over
$\Xx /X$ or $\Xx$), it is HBKQ-cofibrant. One can therefore use it in lemma
\ref{howtocalc}. One obtains the following statement.

\begin{lemma}
\label{calcdonnees}
Suppose that $\Xx$ admits fibre products.
Let $\Bb \subset \Xx /X$ be a sieve generated by a family $\Uu$ covering $X$.
Let $D\rightarrow \ast _{\Bb}$ be the Segal $n$-prestack constructed above.
Then for every Segal $n$-prestack $A$ on $\Xx$ such that all the $A(X)$ are
fibrant, the set of descent data
$$
[\ast _{\Bb}, A|_{\Bb}]_{\Bb ^{\rm gro}}
$$
is computed as the set of morphisms $D\rightarrow A$ modulo the relation
which identifies two morphisms if they are connected by a morphism
$D\times \overline{I}\rightarrow A$.
\end{lemma}
\eop

One can give an explicit description of what a morphism $D\rightarrow A$ is.
One proceeds by induction on the $p$ of $D^p$. If one has a morphism
$D^{p-1}\rightarrow A$, then one obtains for each $\alpha _0,\ldots , \alpha
_p$ the composite morphism
$$
\partial R(p)\rightarrow D^{p-1}(U_{\alpha _0,\ldots ,\alpha _p})\rightarrow
A(U_{\alpha _0,\ldots ,\alpha _p}).
$$
An extension to a morphism $D^p\rightarrow A$ is nothing other than an
extension, for each $\alpha _0,\ldots , \alpha _p$, of the morphism above to a
morphism
$$
R(p)\rightarrow
A(U_{\alpha _0,\ldots ,\alpha _p}).
$$
A morphism $D\rightarrow A$ consists in such a collection of successive
extensions.

In summary: a morphism $\eta : D\rightarrow A$ is a collection of data as
follows: for each $\alpha$, an object $\eta (\alpha )\in A(U_{\alpha})$; for
each $\alpha , \beta$ an equivalence
$$
\eta (\alpha , \beta ): \overline{I}\rightarrow A(U_{\alpha \beta})
$$
between the restrictions of the objects $\eta (\alpha )$ and $\eta (\beta )$;
for each $\alpha , \beta , \gamma$ a morphism
$$
\eta (\alpha , \beta , \gamma ):
\overline{I}^{(2)}\rightarrow A(U_{\alpha \beta\gamma})
$$
which, restricted to the triangle $\partial R(2) \subset \overline{I}^{(2)}$,
is equal to the morphism obtained by restricting $\eta (\alpha ,\beta )$,
$\eta (\beta , \gamma )$ and $\eta (\alpha , \gamma )$; etc.


\subnumero{Comment}
In this description of descent data, we have used in an essential way the
hypothesis that the $A(X)$ are fibrant Segal $n$-categories. It is for this
reason that one does not see $2$-morphisms, $3$-morphisms etc.\ appear in the
data. For example, since $A(U_{\alpha \beta \gamma})$ is fibrant, the datum as
in \cite{Breen} of three $1$-arrows with a $2$-equivalence between $fg$ and
$h$, can be replaced by a genuine morphism from $R(2)$ to
$A(U_{\alpha \beta \gamma})$ (i.e. in some sense an object of
$A(U_{\alpha \beta \gamma})_{2/}$).

This explains why we manage to treat the case of $n$-categories without
getting too entangled in ``cocycle formulas'' but in fact it is a disadvantage
if one seeks an explicit understanding of what is going on, because the notion
of fibrant Segal $n$-category is not very explicitly computable.
\footnote{
From a more optimistic point of view one can nevertheless observe that if $T$
is a topological space (or Kan simplicial set) then $\Pi _{n,Se}(T)$ is
fibrant; so the Segal $n$-prestacks which come from applying $\Pi _{n,Se}$ to
presheaves of topological spaces (or Kan simplicial sets) satisfy the
condition in question.}
In particular, if one does not manage to give a totally ``in terms of
cocycles'' form of the notion of descent data, in the way Breen did for
$2$-stacks in \cite{Breen}, it is because this problem is more difficult than
those one can currently solve. One can formulate a

\begin{problem}
Find a cmc structure analogous to the HBKQ structure, for the (Segal)
$n$-precategories whose fibrant objects are the Segal $n$-categories.
\end{problem}

With an affirmative solution to this problem, one could take up again the
discussion of descent data above: if $D'\rightarrow \ast _{\Bb}$ is a cofibrant
replacement for the structure sought in this problem (which would also give a
cmc structure for the $\Xx^o$-diagrams by the HBKQ method), then one could say
that for every $A$ such that the $A(X)$ are Segal $n$-categories, a descent
datum for $A$ with respect to $\Bb$ is the same thing as a morphism
$D'\rightarrow A$. In that case, one would see (many) $i$-arrows appear for
every $i$, and one would in fact obtain a description of descent data in terms
of cocycles in the classical sense of the term (which would generalise
\cite{Breen}).
